\documentclass[main]{subfiles}
\usepackage[sols]{workbook}
\numbering{ws3-10}

\begin{document}

\ifSubfilesClassLoaded{%
  \chaptermark{\chapterthree}
}{}

\section{Phase Plane Analysis} \label{ws3-10}
\begin{framed}
\centerline{\textbf{Lesson Goals}}

You should be able to:
\begin{itemize}
\item Perform a full phase plane analysis of a given system of first-order differential equations, and use it to determine the long-term behavior of the solution functions under different initial conditions
\item Decide what can and cannot be determined about solution trajectories from the phase plane analysis
\end{itemize}
\end{framed}



\begin{enumerate}  
    
\item  % written by Professor MacMahon

Mrs. Portwick lives on an uninhabited island.  She builds
a small, enclosed garden and decides to import some robins and worms. Let $w(t)$ represent the
number of hundreds of worms on the island at time $t$, and let $r(t)$ be
the number of hundreds of robins in the garden at time $t$.  She
models their interaction as follows:
\begin{align*}
\frac{dr}{dt} & = -0.4r + 0.1wr \\
\frac{dw}{dt} & = 0.3w - 0.3wr
\end{align*}

\begin{enumerate}
\item 
\begin{problem}[0.5in]
Based on the system above, what type of interaction do the robins and worms exhibit? What happens to each species in the absence of the other?
\end{problem}

\begin{solution}
Because the interaction term is positive in the first equation but negative in the second equation, we can determine that the robins and worms have a \fbox{predator-prey} relationship, with robins as the predator and worms as the prey.

In the absence of worms, the robin population will decline exponentially. In the absence of robins, the worm population will increase exponentially.

\end{solution}


\item 
\begin{problem}[\vfill]
Perform a \textit{phase plane analysis} of this system. This means: find and plot all nullclines,
find all equilibrium points, orient the nullclines with direction arrows, and draw direction arrows in each region of the phase plane created by the nullclines. 
\end{problem}

\begin{solution}
We are going to do a phase plane analysis of this system.  This means
that we are going to draw a diagram of the phase plane which shows the
general direction that solution trajectories travel.  Remember that
the phase plane means the $rw$-plane, and we don't get any information
about time $t$ when we look at the phase plane.  By ``general
direction,'' we mean: Does the solution trajectory travel straight up?
Up and right?  Straight to the right?  And so on.

We'll do this in a series of steps.
\begin{itemize}
  \item \textbf{Step 0: Factor the equations.}

  Although it is not strictly necessary to factor the expressions for
  $\frac{dr}{dt}$ and $\frac{dw}{dt}$, it generally makes the rest of
  the phase plane analysis much easier.  I like to factor out a
  constant that makes everything else integers.  In this case, I would
  factor as:
  \begin{align}
  \frac{dr}{dt} & = -0.1 r (4 - w)
  \\
  \frac{dw}{dt} & = 0.3 w (1 - r)
  \end{align}

  \item \textbf{Step 1: Find and plot the $r$-nullclines by figuring out when $\frac{dr}{dt} = 0$.  Draw
  short line segments at these places.}

  First, when $\frac{dr}{dt} = 0$, $r$ is not changing,\footnote{This
  is imprecise.  What we really mean is that the instantaneous rate
  of change of $r$ is 0.} but $w$ may be.  Therefore, we want to
  draw short \textbf{vertical} line segments to represent this.
  (After all, if we were to follow a vertical line segment for a bit,
  $w$ would change, but $r$ would not.)

  It is very easy to figure out when $\frac{dr}{dt} = 0$ if we look at
  the factored equation (1): it happens
  when $r = 0$ or $w = 4$.  So, there are two $\frac{dr}{dt} = 0$
  nullclines.
  As we mentioned already, along each $\frac{dr}{dt} = 0$ nullcline,
  we draw short vertical lines:
  \begin{center}
  \includegraphics[width=2.5in]{ws/images/portwick-1}
  \end{center}
  (Since $r$ and $w$ represent population sizes, we don't care what
  happens when they are negative, so we only draw the first quadrant
  of the $rw$-plane.)

  \item \textbf{Step 2: Find and plot the $w$-nullclines by figuring out when $\frac{dw}{dt} = 0$.  Draw short
  line segments at these places.}

  When $\frac{dw}{dt} = 0$, $w$ is not changing, but $r$ may be.
  Therefore, we want to draw short \textbf{horizontal} line segments
  to represent this.

  Using (2), we see that $\frac{dw}{dt} = 0$ when $w =
  0$ or $r = 1$.
  Along each of these nullclines, we draw short horizontal line
  segments:  
  \begin{center}
  \includegraphics[width=2.5in]{ws/images/portwick-2}
  \end{center}

  \item \textbf{Step 3: Identify the equilibrium points, and draw dots
  there.} 

  The equilibrium points are where $\frac{dr}{dt} = 0$ and
  $\frac{dw}{dt} = 0$.  In other words, these are where the red
  vertical lines and green horizontal lines cross.  You can find these
  just by looking at the sketch so far (they happen at $(0, 0)$ and
  $(1, 4)$).

  We will draw the equilibrium points using blue dots (replacing the
  red and green segments that were there):
  \begin{center}
  \includegraphics[width=2.5in]{ws/images/portwick-3}
  \end{center}

  \item \textbf{Step 4: Orient the $\frac{dr}{dt} = 0$ nullclines.}

  By ``orienting'' the nullclines, we mean that we should draw arrows
  on them to indicate the direction solution trajectories would travel
  with time.

  We already saw in Step 1 that the $\frac{dr}{dt} = 0$ nullclines
  should be indicated by vertical line segments.  Now, we would like
  to turn these line segments into arrows; to decide whether the
  arrows should point up or down, we want to know whether $w$ should
  increase or decrease as time goes on.  That is, we want to know
  whether $\frac{dw}{dt}$ is positive or negative.

  \colorbox{yellow}{ {\LARGE\Pointinghand} \textbf{Key Idea:}
  $\displaystyle \frac{dw}{dt}$ can only change from positive to
  negative if we cross a $\displaystyle \frac{dw}{dt} = 0$
  nullcline.
  }

  In this case, the $\frac{dw}{dt} = 0$ nullclines (shown in green)
  divide the picture into two regions, one to the left of $r = 1$, and
  one to the right of $r = 1$.  Within each region, the sign of
  $\frac{dw}{dt}$ cannot possibly change, so all we need to do is to
  test one point in each region to figure out whether $\frac{dw}{dt}$
  is positive or negative in that region.

  If we test a point like $r = 0, w = 1$ in the left region by
  plugging these values back into the equation $\frac{dw}{dt} = 0.3w
  (1 - r)$, we find that $\frac{dw}{dt} > 0$.  So, we draw up arrows
  along the red nullclines in the left region:
  \begin{center}
  \includegraphics[width=2.5in]{ws/images/portwick-4a}
  \end{center}
  If we test a point in the right region (like $r = 2, w = 1$) by
  plugging these values into the equation $\frac{dw}{dt} = 0.3w (1 -
  r)$, we find that $\frac{dw}{dt} < 0$.  So, we draw down arrows
  along the red nullclines in the right region:
  \begin{center}
  \includegraphics[width=2.5in]{ws/images/portwick-4}
  \end{center}

  \item \textbf{Step 5: Orient the $\frac{dw}{dt} = 0$ nullclines.}

  In Step 2, we found that the $\frac{dw}{dt} = 0$ nullclines were $w
  = 0$ and $r = 1$.  We have drawn them using horizontal line
  segments, and now we want to decide if the arrows should point right
  or left.  To determine this, we want to see if $r$ is increasing or
  decreasing with time, so we look at the sign of $\frac{dr}{dt}$.
  This is completely analogous to what we did in Step 4.

  \colorbox{yellow}{ {\LARGE\Pointinghand} \textbf{Key Idea:}
  $\displaystyle \frac{dr}{dt}$ can only change from positive to
  negative if we cross a $\displaystyle \frac{dr}{dt} = 0$
  nullcline.
  }

  Now, the $\frac{dr}{dt} = 0$ nullclines (shown in red), divide the
  picture into two regions, one above $w = 4$ and one below.  If we
  test one point below $w = 4$, we find that $\frac{dr}{dt} < 0$, so
  the arrows in that region go to the left:
  \begin{center}
  \includegraphics[width=2.5in]{ws/images/portwick-5a}
  \end{center}
  If we test one point above $w = 4$, we find that $\frac{dr}{dt} >
  0$, so the arrows in that region go to the right:
  \begin{center}
  \includegraphics[width=2.5in]{ws/images/portwick-5}
  \end{center}

  \item \textbf{Step 6: Orient each region.}

  Notice that the nullclines have divided the first quadrant of the
  $rw$-plane into 4 regions.  Within a single region, $\frac{dr}{dt}$
  and $\frac{dw}{dt}$ cannot change sign (for instance, if
  $\frac{dr}{dt}$ were to change from positive to negative, it would
  have to be 0, but this can only happen along a $\frac{dr}{dt} = 0$
  nullcline).  So, we can use the orientation of the nullclines that border each region
  to decide the general direction of trajectories within that region (by ``general direction,''
  we mean ``up and left,'' ``up and right,'' ``down and left,'' or
  ``down and right'').

The top-left region is bounded by red nullclines that are oriented upward, and green nullclines that are oriented rightward. So, any solution trajectories passing through that region would have to generally move up and to the right throughout that region. Thus, we will put an arrow pointing in the northeast direction to indicate this fact.

  If you do this for the other 3 regions, you get the following
  diagram:
  \begin{center}
  \includegraphics[width=2.5in]{ws/images/portwick-6}
  \end{center}

  \item \textbf{Step 7: Make sure your sketch is somewhat reasonable.}

  When analyzing systems that model interactions between populations,
  we often think about what would happen if only one of the species
  was present.  For instance, if there are no robins, then the worm
  population should grow exponentially.  (If $r = 0$, then
  (2) becomes $\frac{dw}{dt} = 0.3 w$.)  When we look
  at our diagram, we see that, when $r = 0$ (along the vertical axis),
  the arrows do indeed point up, indicating that the worm population
  increases.

  Similarly, if there are no worms, then the robin population
  decreases exponentially.  In our diagram, the arrows along $w = 0$
  do indeed point left.

  If in doubt, remember that the point of all the arrows is to
  indicate the general direction of the trajectories.  (Up and right?
  Up and left?  Straight up?  And so on.)  You can check this easily
  at any point by plugging in the coordinates of that point to figure
  out whether $\frac{dw}{dt}$ and $\frac{dr}{dt}$ are positive,
  negative, or zero. 
\end{itemize}
\end{solution}



\item 
\begin{problem}
On your phase plane from part (b), sketch trajectories
with the following initial conditions:
\begin{enumerate}[topsep=0pt,itemsep=0pt]
\item $r(0) = 0$, $w(0) = 2$.
\item $r(0) = 2$, $w(0) = 0$.
\end{enumerate}
\end{problem}


\begin{solution}
The first trajectory is in red; the second in blue:
\begin{center}
\includegraphics[width=3in]{ws/images/portwick-axis-trajectories}
\end{center}
\end{solution}

\end{enumerate}


\pspace{\newpage}


\item 
Mrs. Portwick's nephew Edwin comes to visit, sees her work, and says
that he believes that the second differential equation should be
$\displaystyle \frac{dw}{dt} = 0.3w - 0.03 w^2 - 0.3wr$.

\begin{enumerate}
  \item
  \begin{problem}[0.5in]
  What could his reasoning be?
  \end{problem}

  \begin{solution}
  In the original model, if there were no robins ($r = 0$), then
  $\frac{dw}{dt} = 0.3 w$, so the worm population would always grow at
  a rate proportional to its size (which we know results in
  exponential growth).

  In Edwin's suggested model, if there are no robins, then
  $\frac{dw}{dt} = 0.3 w - 0.03 w^2 = 0.03 w (10 - w)$.  You should
  recognize this as a logistic model with carrying capacity $w = 10$
  ($1000$ worms since $w(t)$ represents hundreds of worms).  Edwin's
  reasoning is that exponential growth is unrealistic, and the garden
  can only support a limited number of worms ($1000$).
  \end{solution}

  \item

  \note{Students tend to have a lot of trouble with the nullcline $10 -
  w - 10r = 0$.  Some common errors:
  \begin{itemize}[label={--}] 
    % label avoids Font shape `OMS/cmss/m/n' undefined error for bullets
    % in framebox? JA
    \item For whatever reason, students often don't recognize $10 - w -
    10r = 0$ as describing a line.  It may help to talk about graphing
    $10 - y - 10x = 0$ instead.
    \item Some students see $10 - w - 10r = 0$ and start thinking of
    individual points that make this true, like $r = 1, w = 0$ or $w =
    10$, $r = 0$.  Emphasize that we need to find \textbf{all} points
    where this is 0.
    \item It is fairly common for students to try to draw the little
    line segments / arrows \textbf{along} the nullcline, rather than
    as horizontal things.
  \end{itemize}
  I recommend just letting them make these errors and then talking
  about why they are wrong after the mistakes have been made.
  (Warning students in advance doesn't seem to actually prevent them
  from making the errors.)}

  \begin{problem}[\vfill]
  Perform a phase plane analysis on the following new system:
  \begin{align*}
\frac{dr}{dt} & = -0.4r + 0.1wr \\
\frac{dw}{dt} &= 0.3w - 0.03 w^2 - 0.3wr
\end{align*}
  \end{problem}

  \begin{solution}
  We follow exactly the same steps to analyze this system.
  \begin{itemize}
    \item \textbf{Step 0: Factor the equations.}
    \begin{align}
    \frac{dr}{dt} & = -0.1 r (4 - w)
    \\
    \frac{dw}{dt} & = 0.03 w (10  - w - 10r)
    \end{align}

    \item \textbf{Step 1: Figure out when $\frac{dr}{dt} = 0$ to find and plot the $r$-nullclines.  Draw
    short (vertical) line segments at these places.}

    By (3), $\frac{dr}{dt} = 0$ when $r = 0$
    or $w = 4$.  For future reference, let's call these 2 nullclines
    r1 and r2:
    \begin{equation*}
    r = 0 \tag{r1}
    \end{equation*}
    \begin{equation*}
    w = 4 \tag{r2}
    \end{equation*}
    \begin{center}
    \includegraphics[width=2.5in]{ws/images/portwick-logistic-1}
    \end{center}

    \item \textbf{Step 2: Figure out when $\frac{dw}{dt} = 0$ to find and plot the $w$-nullclines.  Draw
    short (horizontal) line segments at these places.}

    By (4), $\frac{dw}{dt} = 0$ when $w = 0$
    or $10 - w - 10r = 0$.  So, we again have 2 nullclines.  The
    second one is still a line (the line $w = 10 - 10r$); it's just
    not parallel to either axis.

    \settowidth{\dwidth}{\Large\danger}
    {\Large\danger}
    \begin{minipage}[t]{\linewidth-\dwidth-2\fboxsep}
    A common mistake is to look at $10 - w - 10r = 0$ and think
    something like, ``Oh, that will be true when $w = 0$ and $r =
    1$,'' and leave things at that.  The $\frac{dw}{dt} = 0$
    nullclines represent \textbf{ALL} points where $\frac{dw}{dt} =
    0$, so just thinking of a few values that make the equation true
    is not enough.  Instead, think about what shape the equation
    describes.  (It's easiest if you can solve one variable for the
    other, like how we rewrote the equation $10 - w - 10r = 0$ as $w
    = 10 - 10r$.)
    \end{minipage}

    Again, let's name them for future
    reference:
    \begin{equation*}
    w = 0 \tag{w1}
    \end{equation*}
    \begin{equation*}
    10 - w - 10r = 0 \tag{w2}
    \end{equation*}
    \begin{center}
    \includegraphics[width=2.5in]{ws/images/portwick-logistic-2}
    \end{center}

    \item \textbf{Step 3: Identify the equilibrium points, and draw dots
    there.} 

    By looking at where the red vertical lines and green horizontal
    lines cross, we see that there are 3 equilibrium points, $(0, 0)$,
    $(0, 10)$, and $\left( \frac{3}{5}, 4 \right)$.

    Alternatively, we can look at the equations we've found for the
    nullclines.  An equilibrium point is the intersection of a
    $\frac{dr}{dt} = 0$ nullcline (r1 or r2) and a $\frac{dw}{dt} = 0$
    nullcline (w1 or w2).  So, there are 4 possibilities to consider:
    \begin{itemize}
      \item ({r1}) and
      ({w1}): $r = 0$ and $w = 0$, which gives
      us the equilibrium point $(0, 0)$.
      \item ({r1}) and
      ({w2}): $r = 0$ and $10 - w - 10r = 0$.
      When $r = 0$, the second equation is just $10 - w = 0$, so this
      gives the equilibrium point $(0, 10)$.
      \item ({r2}) and
      ({w1}): $w = 4$ and $w = 0$, but these
      can't both be true at the same time.
      \item ({r2}) and
      ({w2}): $w = 4$ and $10 - w - 10 r = 0$.
      When $w = 4$, the second equation is just $6 - 10r = 0$, so this
      gives us the equilibrium point $\left( \frac{3}{5}, 4 \right)$.
    \end{itemize}

    So, here is our diagram with the 3 equilibrium points.
    \begin{center}
    \includegraphics[width=2.5in]{ws/images/portwick-logistic-3}
    \end{center}

    \item \textbf{Step 4: Orient the $\frac{dr}{dt} = 0$ nullclines by
    looking at the sign of $\frac{dw}{dt}$.}

    Remembering that $\frac{dw}{dt}$ can only change sign when we
    cross a $\frac{dw}{dt} = 0$ nullcline (marked in green), we see
    that there are two regions to look at, the one to the left of the
    nullcline ({w2}) and the one to the right.

    We can test one point in each region, but let's instead look at
    what happens along the nullcline $r = 0$.  (I like to do this
    because it's one of the things we look at anyway when interpreting
    models.)  When $r = 0$ (there are no robins), equation
    (4) becomes $\frac{dw}{dt} = 0.03 w (10 -
    w)$.  This is an autonomous differential equation, and we can use
    qualitative analysis to draw a family of representative solutions: 
    \begin{center}
    \includegraphics[width=2in]{ws/images/portwick-logistic-worms}
    \end{center}
    From this, we see that, when $r = 0$, $\frac{dw}{dt} > 0$ for $w$
    between 0 and $10$ and $\frac{dw}{dt} < 0$ for $w > 10$.  This
    matches our answer from {{\#3}(a)}
    perfectly: without any robins, the worms grow logistically with
    carrying capacity 10.  This allows us to fill in the arrows along
    the line $r = 0$:
    \begin{center}
    \includegraphics[width=2.5in]{ws/images/portwick-logistic-4a}
    \end{center}
    Now, notice that we have already filled in some arrows in each
    of the two regions we were talking about.  In the region to the
    left of nullcline ({w1}), we have drawn some
    up arrows, so we can fill the rest of that region in with up
    arrows as well.  In the region to the right of nullcline
    ({w2}), we have drawn some down arrows, so
    we can fill the rest of that region in with down arrows as well.
    So, we have the following diagram:
    \begin{center}
    \includegraphics[width=2.5in]{ws/images/portwick-logistic-4}
    \end{center}

    \item \textbf{Step 5: Orient the $\frac{dw}{dt} = 0$ nullclines by
    looking at the sign of $\frac{dr}{dt}$.}


    Let's try the same strategy here.  We really need to look only at
    two regions, the one above $w = 4$ and the one below $w = 4$.  But
    let's start by thinking about what happens along the $r$-axis, when
    $w = 0$.  When $w = 0$, equation (3)
    becomes $\frac{dr}{dt} = -0.4 r$, which we recognize as
    exponential decay; that is, the robin population is always
    falling.  So, all of the arrows along $w = 0$ should go to the
    left.  From this, we can also fill in the rest of the arrows under
    $w = 4$, to get this diagram:
    \begin{center}
    \includegraphics[width=2.5in]{ws/images/portwick-logistic-5a}
    \end{center}
    This time, we have no information about the upper region, so let's
    just test a point there, like $r = 1$, $w = 5$.  If we plug this
    into equation (3), we find that
    $\frac{dr}{dt} > 0$ here, so we need to draw right arrows in this
    region.  So, we end up with the following picture:
    \begin{center}
    \includegraphics[width=2.5in]{ws/images/portwick-logistic-5}
    \end{center}

    \item \textbf{Step 6: Orient each region.}

    Finally, we use the orientation of the boundaries for each of the resulting regions to
    get: 
    \begin{center}
    \includegraphics[width=2.5in]{ws/images/portwick-logistic-6}
    \end{center}

  \end{itemize}
  \end{solution}
  
  

  
  \item 
\begin{problem}
On your phase plane analysis from part (b), sketch
trajectories with the following initial conditions:
\begin{enumerate}[topsep=0pt,itemsep=0pt]
\item $r(0) = 0$, $w(0) = 2$.
\item $r(0) = 2$, $w(0) = 0$.
\item $r(0) = 0$, $w(0) = 12$.
\end{enumerate}
\end{problem}

\begin{solution}
The first trajectory is in red, the second is in green, and the third
is in blue:
\begin{center}
\includegraphics[width=3in]{ws/images/portwick-logistic-axis-trajectories}
\end{center}
\end{solution}
  
\end{enumerate}

  
\pspace{\newpage}








For the following questions, assume that we are modeling Mrs. Portwick's
garden using the model in {\#1}. Here is the phase plane for your convenience:

  \begin{center}
  \includegraphics[width=2.25in]{ws/images/portwick-6}
  \end{center}

\item 

\textit{It turns out that, if the robin and worm populations are both positive, the solution trajectories of the system in
{\#1} are closed loops.  (This is \textbf{NOT} something we
can tell from the phase plane analysis.)}

Suppose that, at some time $t = 0$, Mrs. Portwick's
garden has 200 robins ($r = 2$) and 400 worms ($w = 4$).
\begin{enumerate}
  \item
  \begin{problem}[0.5in]
  Plot the solution trajectory with the given initial condition on the phase plane above. Describe in words what happens to the populations over time.  
  \end{problem}

  \begin{solution}
    This is the blue trajectory in the picture below.
    
  \begin{center}
    \includegraphics[width=2.25in]{ws/images/portwick-blue.eps}
  \end{center}
  
  At $t = 0$, we
  are at the right-most point on this trajectory.  As we follow this
  trajectory, the robin and worm populations both decrease.  (There
  are too many robins compared to worms.  The robins are eating so
  many worms that the worm population is decreasing, but there is not
  really enough food for the robins, so the robin population is also
  decreasing.)  When we get to the lowest point of the trajectory ($r
  = 1$, $w \approx 1.9$; let's say this happens at time $T_1$), the
  worm population starts to increase, but the robin population
  continues to decrease.  (There are few enough robins that the worm
  population can grow, but the robins still don't have enough food, so
  the robin population continues to decline.)  This continues until we
  reach the left-most point of the trajectory ($r \approx 0.4$, $w =
  4$; let's say this happens at time $T_2$).  At this point, the robin
  population starts to increase again, while the worm population
  continues to increase.  (There are plenty of worms, so the robin
  population starts to increase; there are more worms than the robins
  can eat, so the worm population continues to increase.)  When we
  reach the upper-most point of the trajectory ($r = 1$, $w \approx
  7.4$; let's say this happens at time $T_3$), the worm population
  starts to decrease, while the robin population continues to
  increase.  (The robins have plenty of food, so the robin population
  continues to increase.  But they are eating so many worms that the
  worm population decreases.)  Finally, we get back to $r = 2, w = 4$
  (let's say this happens at time $T_4$), and the whole cycle starts
  over.
  \end{solution}

  \item
  \begin{problem}[0.25in]
  On the following set of axes, draw possible graphs of $r(t)$ and
  $w(t)$. 
  \end{problem}

  \begin{problemonly}
  \includegraphics[height=1.75in]{ws/images/phase-portrait-time-graphs-axes-first-quadrant} 
  \end{problemonly}

  \begin{solution}
  Based on the description in words, the graphs might look something
  like this, with $r(t)$ in red and $w(t)$ in blue: 
  \begin{center}
  \includegraphics[width=2in]{ws/images/portwick-time-graphs}
  \end{center}
  The main points to note are: at $T_1$, $r = 1$ and $w \approx 1.9$;
  at $T_2$, $r \approx 0.4$, $w = 4$; at $T_3$, $r = 1, w \approx
  7.4$; at $T_4$, $r = 2, w = 4$.  The graphs are periodic with period
  $T_4$.

  Remember, though, that the phase plane analysis really
  doesn't tell us how long anything takes (in other words, we don't
  have any idea what times $T_1$, $T_2$, $T_3$, and $T_4$ are).  So,
  you would be just as correct if you drew your graphs with $T_1$,
  $T_2$, $T_3$, and $T_4$ equally spaced or with some other spacing.
  (To really determine the proper spacing for $T_1$, $T_2$, $T_3$, and
  $T_4$, we would need to go back to the original system; for this
  particular system, this is not something we'll do in Math 1b.
  If you'd like to learn more about solutions to these types of problems,
  we recommend taking a course in differential equations or dynamical systems.)
  \end{solution}

\end{enumerate}

\item 

\note{This one is kind of interesting and perhaps a little
counter-intuitive.  It has interesting applications in all sorts of
pest control --- for instance, if you spray mosquitoes at the wrong
time, then you might end up with more of them at peak times.}

\begin{problem}[1in]
One day, Mrs. Portwick observes that there are 150 robins and 300 worms.
Mrs. Portwick feels that the garden has too many worms, so she removes
50 of them to use as fishing bait.  What effect does
this have on the worm population in the long run?
\end{problem}

\begin{solution}
If Mrs. Portwick had left things as they were, the system would have
been on the red trajectory, which goes through $(1.5, 3)$.
\begin{center}
\includegraphics[width=2in]{ws/images/portwick-modify-trajectories}
\end{center}
Her meddling puts the system on the blue trajectory, which goes
through $(1.5, 2.5)$.  In the long run, the worm population will
fluctuate more than it did before; at its low points, there will be
fewer worms than there would have been had she not meddled.  However,
at the worm population's high points, there will be \textit{more} worms than
there would have been!
\end{solution}


  \pspace{\newpage}


The following problems are here to give you extra practice with phase plane analysis. In these examples, the nullclines are curves instead of straight lines. But the process remains the same!

\item \label{trapping-region}
These questions concern the system:
$$
  \left\{ \begin{aligned}
  \frac{dx}{dt} & = y^2 - x \\
  \frac{dy}{dt} & = x^2 - y
  \end{aligned} \right.
$$

  \begin{enumerate}
    \item
    \begin{problem}[3in]
      Perform a phase plane analysis of this system.
    \end{problem}

    \begin{solution}
    There is no need to factor the right-hand sides of these differential
    equations.  Setting $dx/dt = 0$, we obtain the $x$-nullcline $x = y^{2}$,
    which is a parabola.  Setting $dy/dt = 0$, we obtain the $y$-nullcline
    $y = x^2$, which is also a parabola.  In Panel (a) of the following figure,
    the $x$-nullcline is indicated in blue, and the $y$-nullcline is indicated
    in red.  {\bf WARNING:}  In this figure, the nullclines are rendered as
    continuous curves, but keep in mind that they are NOT solution trajectories!
    Any solution trajectory that crosses the blue curve must have a vertical
    tangent at the crossing point, and any solution trajectory that crosses the
    red curve must have a horizontal tangent at the crossing point.
    There are two equilibria, $(x,y) = (0,0)$ and $(x,y) = (1,1)$,
    indicated by bold dots in the figure.  

    Now, let's orient the nullclines, starting with the $x$-nullcline.  If
    $x = y^2$, we know that $dx/dt = 0$.  What can we say about $dy/dt$ along
    this nullcline?  If we substitute $x = y^2$ into the equation for $dy/dt$,
    we find that
    $$
	\frac{dy}{dt} = y^4 - y \; = \; y (y^3 - 1)
    $$  
    at each point along the [blue] $x$-nullcline.  Along the portion of the
    $x$-nullcline for which $y < 0$, we see that $dy/dt > 0$.  In Panel (b)
    of the figure, we've indicated this by sketching upward-pointing arrows
    along that portion of the $x$-nullcline.  Along the portion of the 
    $x$-nullcline for which $0 < y < 1$, we have $dy/dt < 0$.  We've indicated
    this by drawing downward pointing arrows along this portion of the blue
    nullcline.  Finally, along the portion of the $x$-nullcline for which 
    $y > 1$, we have $dy/dt > 0$, as indicated by an upward pointing arrow
    in the figure.

    Orienting the $y$-nullcline is similar.  If
    $y = x^2$, we know that $dy/dt = 0$ so let's focus on the sign of $dx/dt$.
    If we substitute $y = x^2$ into the equation for $dx/dt$, we find that
    $$
	\frac{dx}{dt} = x^4 - x \; = \; x (x^3 - 1)
    $$  
    at each point along the [red] $y$-nullcline.  From here, the reasoning is
    identical to how we proceeded in orienting the $x$-nullcline, and the
    results are indicated in Panel (c) of the figure.

    Finally, in Panel (d) of the figure, we've drawn a few diagonal arrows
    to indicate the general direction of the flow in the various regions
    into which the nullclines have partitioned the phase plane.  


    \begin{center}
      \includegraphics[width=5in]{ws/images/7a.eps}
    \end{center}

    \end{solution}

    \item
    \begin{problem}[0.5in]
      There is a bounded region in the phase plane that is enclosed by two
      nullclines.  If you were to start from initial conditions within
      that enclosed region, does your phase plane analysis from part (a)
      allow you to draw conclusions about the long-term behavior of $x$ and $y$?
    \end{problem}

    \begin{solution}
    Notice
    that for initial conditions inside the bounded region enclosed by the
    two nullclines, we'll remain {\em trapped} inside that region as $t$
    increases!  Apparently, for initial conditions in that region, we would
    expect to approach the $(0,0)$ equilibrium in the long run (which is 
    indeed what happens).  In fact, the phase plane diagram suggests that
    $(0,0)$ is a stable equilibrium:  For every initial condition that's
    appropriately close to $(0,0)$, we'd expect the corresponding solution
    trajectory to approach $(0,0)$ as $t$ increases.  On the other hand,
    it appears that $(1,1)$ is unstable.  If we were to choose 
    $(x,y) = (1.01, 1.01)$ as initial conditions, the phase plane diagram
    suggests that the resulting trajectory would head ``northeast", away
    from the equilibrium.
    \end{solution}



  \end{enumerate}


\item 
\begin{problem}
Try repeating the previous exercise for the system
$$
  \left\{ \begin{aligned}
  \frac{dx}{dt} & = y^2 - x \\
  \frac{dy}{dt} & = y - x^2.
  \end{aligned} \right.
$$
{\bf Remark:}  This system has the same nullclines as in the preceding exercise,
but the phase plane analysis yields very different results!
\end{problem}


    \begin{solution}
    As stated in the remark, this system has the same nullclines as the
    in the preceding exercise:  The $x$-nullcline is $x = y^2$, the blue
    curve in Panel (a) of the following figure.  The $y$-nullcline is
    $y = x^2$, indicated in red in Panel (a) of the following figure.
    
    The $y$-nullcline has precisely the same orientation as in the 
    previous exercise.  Indeed, if $y = x^2$, then $dx/dt = x(x^3 - 1)$ 
    as before.  However, all of the arrows along the $x$-nullcline are
    flipped from how they were in the previous problem.  To see why,
    suppose $x = y^2$, so that we're standing on the $x$-nullcline.  We
    know that $dx/dt = 0$, so let's check the signs of $dy/dt$.  Along
    the curve $x = y^2$, we have $dy/dt = y - y^4 = y(1-y^3)$.  Because
    this equation has precisely the opposite sign as its counterpart in
    the preceding exercise, all of the arrows on the $x$-nullcline are
    flipped.  All of this is summarized in the remaining panels in the
    figure below.  

    This time, the phase plane diagram suggests that solutions near
    $(1,1)$ exhibit oscillatory behavior.  Trajectories near that
    equilibrium may form closed loops surrounding $(1,1)$ (indicative of
    periodic solutions).  Then again, they may spiral inwards towards
    $(1,1)$, or outward away from $(1,1)$.  In Math 1b, we won't have
    the tools to determine whether the solutions near $(1,1)$ really 
    are periodic\footnote{It turns out that they are.}, but at least we
    were able to detect oscillations.

    Take a moment to examine what happens in the vicinity of the other
    equilibrium, $(0,0)$.  In the previous exercise, it seemed pretty
    clear that $(0,0)$ was a stable, attracting equilibrium.  This
    time, it's not as clear whether $(0,0)$ is stable\footnote{It turns 
    out to be unstable.} or not, and that's not something that our 
    phase plane analysis alone could predict.
    
    
    \begin{center}
      \includegraphics[width=5in]{ws/images/8a.eps}
    \end{center}

    \end{solution}




  
\end{enumerate}



 \end{document}
 


\end{document}